\subsection{The Value of Adaptivity in a Structured Model}\label{app:structured}
The constants $V$ and $W_\A(P)$ are for an analyst who assumes nothing about
the law of the unasked features. Here the law of $(X,Y)$ is known to lie in a
model, and we ask what remains of the gap between pre-committed and branching
questions. The answer is the content of Theorem~\ref{thm:structuredfull}: the
simulation bound survives in every model; when the features are independent given the label a design matters only through its inclusion probabilities, and the gap is at most a factor two; in a model with interactions of
order at most $k$ the exponential gap becomes polynomial in the cap at
product laws and near them, and it disappears for the address contrast when
$k=2$; but parsimony alone does not help, because pairwise laws with strong
couplings can encode the address structure and bring the exponential gap
back.

\paragraph{Constants in a model.}
Let $P$ have full support and let $\mathcal T\subseteq L^2_0(P)$ be a linear
space of mean-zero functions of $(X,Y)$, the tangent space at $P$ of a smooth
parametric model $\mathcal M\ni P$. For an interview $\iota$ and $h\in L^2(P)$
write $J_\iota(h)=\E_P\big[\E_P[h\mid\ell_\iota(X),Y]^2\big]$, the Fisher
information that the records of $\iota$ carry about the direction $h$, and
for a design $q$ on a menu $\A$ let $J_q=\sum_\iota q_\iota J_\iota$. Define
\begin{equation}
 W^{\mathcal T}_\A(P)=\inf_q\ \sup_{h\in\mathcal T}\ \frac{\langle t-\Delta(P),h\rangle_P^2}{J_q(h)},\label{eq:WT}
\end{equation}
with $0/0=0$ and the value $+\infty$ if $\langle t-\Delta,h\rangle_P\ne0=J_q(h)$
for some $h\in\mathcal T$. If $s=(s_1,\ldots,s_p)$ is a basis of
$\mathcal T$, $v=\E_P[(t-\Delta)s]$ and
$I_\iota=\E_P\big[\E_P[s\mid\ell_\iota,Y]\E_P[s\mid\ell_\iota,Y]^{\mathsf T}\big]$,
then $W^{\mathcal T}_\A(P)=\min_qv^{\mathsf T}(\sum_\iota q_\iota I_\iota)^\dagger v$,
with the value $+\infty$ when $v$ is outside the range of the matrix, a
$c$-optimal design problem. For $\mathcal T=L^2_0(P)$ this is the constant
$W_\A(P)$ of \eqref{eq:mainW} (proof of Proposition~\ref{prop:local}), and
$\mathcal T\subseteq\mathcal T'$ implies
$W^{\mathcal T}_\A(P)\le W^{\mathcal T'}_\A(P)\le W_\A(P)$: a model can only
lower the constant of a menu, and the question is whether it lowers the
constant of fixed forms by more than that of branching.

\begin{proposition}[Local minimax risk in a model]\label{prop:localmodel}
Let $\{P_\theta:\theta\in\Theta\subseteq\mathbb R^p\}$ be a model on the finite
space $\X\times\{-1,1\}$ with $\theta_0$ in the interior of $\Theta$,
$P_{\theta_0}=P$ of full support,
$\theta\mapsto P_\theta$ continuously differentiable with a Jacobian of rank
$p$ at $\theta_0$, and let $\mathcal T$ be the span of its scores at
$\theta_0$. For every sequence of protocols that choose each record's
interview from $\A$ using the previous records, and every sequence of
estimators,
\[
 \lim_{c\to\infty}\liminf_{n\to\infty}\ \sup_{\|\theta-\theta_0\|\le c/\sqrt n}
     n\,\E_\theta\big[(\widehat\Delta-\Delta(P_\theta))^2\big]\ \ge\ W^{\mathcal T}_\A(P).
\]
\end{proposition}
\begin{proof}
This is the converse half of the proof of Proposition~\ref{prop:local} with
the tangent coordinates $P+U\theta$ of the simplex replaced by the model's
own parameter, and with two changes. The cube of radius $c/\sqrt{pn}$ around
$\theta_0$ that carries the smooth product prior lies in $\Theta$ for large
$n$ because $\theta_0$ is interior. On this cube the information matrix of
one record of interview $\iota$ is $I_\iota(\theta)=I_\iota+o(1)$ and the
gradient of $\theta\mapsto\Delta(P_\theta)$, which is no longer linear, is
$v+o(1)$, uniformly, by continuous differentiability. The multivariate
van Trees inequality then bounds the normalized Bayes risk below by
$(v+o(1))^{\mathsf T}(\sum_\iota\bar q_\iota I_\iota+C_0/c^2+\epsilon_nI)^{-1}(v+o(1))$
for some $\bar q$ in the simplex; the matrix is bounded below by $C_0/c^2$
uniformly in $\bar q$, so the $o(1)$ terms vanish in the limit for fixed $c$.
Minimizing over $\bar q$ and letting $n\to\infty$ and then $c\to\infty$ gives
$\min_qv^{\mathsf T}(\sum_\iota q_\iota I_\iota)^\dagger v$.
\end{proof}
Like $W_\A(P)$ in Proposition~\ref{prop:local}, the constant is local, with
the center $P$ available to the procedure: when $\sum_\iota q_\iota I_\iota$
is nonsingular at a minimizing design, the maximum-likelihood estimator of
$\Delta(P_\theta)$ under that design attains it for regular models such as
the log-linear models below. A minimizing design may have singular
information, and then $\Delta(P_\theta)$ need not be identified from its
records away from $P$; this happens for the address contrast under the
pairwise model below, whose optimal designs read only data bits. The value
is then approached, up to a factor $1/(1-\epsilon)$, by mixing the design with
a mass $\epsilon$ of interviews that identify the model. We use the lower
bound, the comparison of menus, and the simulations of
Appendix~\ref{app:modelexp}.

\paragraph{Models of bounded interaction order.}
Number the label as coordinate $m+1$. The \emph{order-$k$ model}
$\mathcal M_k$ is the hierarchical log-linear model
$\log p(x,y)=\sum_{A:\,|A|\le k}\theta_A\chi_A(x,y)$ over
$A\subseteq[m+1]$; $k=2$ has main effects and pairwise interactions and
$k=m+1$ is saturated. Its tangent space at any of its laws is
$\mathcal T_k=\mathrm{span}\{\chi_A-\E_P\chi_A:1\le|A|\le k\}$, and we write
$W^{(k)}_\A(P)$ for \eqref{eq:WT} with $\mathcal T=\mathcal T_k$. A
\emph{product law} makes $X_1,\ldots,X_m,Y$ independent; it lies in every
$\mathcal M_k$.

\begin{theorem}[Adaptivity in a model, full statement of Theorem~\ref{thm:structured}]\label{thm:structuredfull}
Let $1\le b\le m$ and let $d$ be a contrast of degree at most $b$.

\textup{(a) Every model.} For every depth-$b$ tree $\tau$ with padded path
sets $\mathcal S_\tau$ and every law $P$, the uniform design on
$\mathcal S_\tau$ satisfies
\begin{equation}
 \frac1{|\mathcal S_\tau|}\sum_{S\in\mathcal S_\tau}J_S(h)\ \ge\ \frac{J_\tau(h)}{|\mathcal S_\tau|}
 \qquad\text{for every }h\in L^2(P),\label{eq:infodom}
\end{equation}
where $J_S$ is the information of the form $S$. Consequently, for every
tangent space $\mathcal T$ and $1\le r\le b$,
\[
 W^{\mathcal T}_{\rm branch}(P)\ \le\ W^{\mathcal T}_r(P)\ \le\ \min\Big\{2^{b-r},\binom{m-r}{b-r}\Big\}\,W^{\mathcal T}_{\rm branch}(P).
\]

\textup{(b) Features independent given the label.} Let $\mathcal M_{\rm ci}$
be the model of all full-support laws under which $X_1,\ldots,X_m$ are
independent given $Y$, with tangent space $\mathcal T_{\rm ci}$ and constants
$W^{\rm ci}_\A$; for designs whose interview probabilities $q_\iota(y)$ may depend on the record's label, $W^{\rm ci,lab}_\A$ is \eqref{eq:WT} with $J_q(h)=\sum_yP(y)\sum_\iota q_\iota(y)\,\E[\E[h\mid\ell_\iota(X),Y]^2\mid Y=y]$ (Proposition~\ref{prop:localmodel} applies to such designs with the menu of pairs of interviews, one for each label). For $P\in\mathcal M_{\rm ci}$ let $\phi^y_j$ be $X_j$
standardized under $P(\cdot\mid Y=y)$,
$\hat d^{\,y}_j=\E_P[d(X)\phi^y_j(X_j)\mid Y=y]$, and
\[
 M_b(c)=\min\Big\{\sum_{j=1}^m\frac{c_j^2}{\pi_j}:\ 0\le\pi_j\le1,\ \sum_{j=1}^m\pi_j\le b\Big\}\qquad(0/0=0).
\]
Then, for every contrast $d$ and every law $P\in\mathcal M_{\rm ci}$:
\begin{enumerate}
\item[\textup{(i)}] if the design may depend on the record's label, fixed
forms, $r$-round interviews for every $r$ and branching interviews have the
same constant,
$W^{\rm ci,lab}(P)=\Var_P(\E_P[t\mid Y])+\sum_yP(y)\,M_b(\hat d^{\,y})$;
\item[\textup{(ii)}] for designs chosen without the label,
$W^{\rm ci,lab}(P)\le W^{\rm ci}_{\rm branch}(P)\le W^{\rm ci}_{r}(P)\le W^{\rm ci}_{\rm fixed}(P)=\Var_P(\E_P[t\mid Y])+M_b(\bar c)\le2\,W^{\rm ci,lab}(P)$,
where $\bar c_j^{\,2}=\sum_yP(y)(\hat d^{\,y}_j)^2$;
\item[\textup{(iii)}] if $P$ is a product law with $P(Y=1)=1/2$, then also in
the pairwise model
$W^{(2)}_{\rm fixed}(P)=W^{(2)}_{r}(P)=W^{(2)}_{\rm branch}(P)=(\E_Pd)^2+M_b(\hat d)$
for every $r$, with $\hat d_j=\E_P[d(X)\phi_j(X_j)]$.
\end{enumerate}
The minimum is attained by $\pi_j=\min\{1,|c_j|/\lambda\}$ for the $\lambda$
that makes the sum equal to $b$ (and by $\pi_j=\mathbf 1\{c_j\ne0\}$ if at most $b$ of the
$c_j$ are nonzero); it equals $(\sum_j|c_j|)^2/b$ when
$b\max_j|c_j|\le\sum_j|c_j|$.

\textup{(c) Bounded interaction order at product laws.} Let $P$ be a product
law whose coordinates and label take each value with probability at least
$\lambda\in(0,1/2]$, and let $2\le k\le b+1$. For every depth-$b$ tree
$\tau$, the fixed design $q^\tau_S=P(S_\tau(X)=S)$ satisfies
$\sum_Sq^\tau_SJ_S(h)\ge J_\tau(h)/\nu_k$ for every $h\in\mathcal T_k$,
with
\[
 \nu_k=\nu_k(b,\lambda)=\sum_{j=1}^{k}\binom{b+1}{j}\Big(\frac{1-\lambda}{\lambda}\Big)^{j},
\]
and therefore
\[
 W^{(k)}_{\rm fixed}(P)\ \le\ \min\{K_{m,b},\ \nu_k(b,\lambda)\}\ W^{(k)}_{\rm branch}(P).
\]
At uniform laws $\nu_k=\sum_{j=1}^k\binom{b+1}{j}\le(b+1)^k$.
The bound is stable: if $k\le b$ and $P_0$ is such a product law, then for
every $\epsilon>0$ there is a neighbourhood of $P_0$ in $\mathcal M_k$ on
which $W^{(k)}_{\rm fixed}(P)\le(1+\epsilon)\,\nu_k(b,\lambda)\,W^{(k)}_{\rm branch}(P)$.

\textup{(d) The address contrast.} Let $d$ be the address contrast with
$\kappa=b-1$ address bits and $D=2^\kappa$ data bits, and $P$ the uniform law
with an independent uniform label, so that $\sigma_P^2=1$ and
$W_{\rm fixed}(P)/W_{\rm branch}(P)=2^{b-1}$. Write
$N(s,j)=\sum_{i=0}^{j}\binom si$. Then
\[
 W^{(2)}_{\rm fixed}(P)=W^{(2)}_{\rm branch}(P)=\frac1b,
\]
and for $3\le k\le b+1$,
\[
 W^{(k)}_{\rm branch}(P)\le1,\qquad
 W^{(k)}_{\rm fixed}(P)\ \ge\ L_k(b)=\frac{N(b-1,k-2)^2}{\max_{0\le s\le b-1}(b-s)\,N(s,k-2)}\ \ge\ \frac1b\binom{b-1}{k-2}.
\]
Here $L_3(b)=b^2/\lfloor(b+1)^2/4\rfloor<4$ and $L_{b+1}(b)=2^{b-1}$. For
$k=3$, moreover, $W^{(3)}_{\rm fixed}(P)<4$ for every $b\ge2$, and already
two-round interviews give
\[
 W^{(3)}_{\rm branch}(P)\ \le\ U_3(b)=\min_{\substack{s+r=b,\ s,r\ge1\\ r\le2^{\kappa-s}}}
   \Big(\frac{\kappa}{s\sqrt r}+\frac{\kappa-s}{s\sqrt b}\Big)^2=\frac{9+o(1)}{b},
\]
so that $W^{(3)}_{\rm fixed}(P)/W^{(3)}_{\rm branch}(P)\ge L_3(b)/U_3(b)=(4/9+o(1))\,b$.

\textup{(e) Strong dependence.} Let $b\ge3$, $\kappa=b-1$ and $D=2^\kappa$.
Add to the $\kappa+D$ coordinates of the address contrast $2D$ auxiliary
binary features $s,p\in\{0,1\}^D$, which $d$ does not depend on; write $z_j$
for the indicator of $Z_j=1$. (Coding a feature by $\{0,1\}$ is an affine
change of the $\pm1$ coordinate and leaves every $\mathcal M_k$ unchanged.)
For every $\epsilon>0$ there is a full-support law $P$ in the pairwise model
$\mathcal M_2$ on these $m=b-1+3\cdot2^{b-1}$ features and the label, under
which, with probability at least $1-\epsilon$, $s$ is the one-hot code of the
address and $p_j=s_jz_j$ for all $j$, such that for every $k\ge2$,
\[
 \frac{W^{(k)}_{\rm fixed}(P)}{W^{(k)}_{\rm branch}(P)}\ \ge\ (1-\epsilon)\,\frac{2^{b-1}}{b}.
\]
For $b=2$ the same construction gives only the trivial bound one.
\end{theorem}

Part (a) is the model-level form of Theorem~\ref{thm:simulation}: with
probability $1/|\mathcal S_\tau|$ a record of the uniform path-set design
contains the transcript that $\tau$ would have produced, whatever the analyst
intends to do with it, so the factor $2^{b-1}$ bounds the value of branching
for maximum likelihood in any model, for Bayesian analysis and for
imputation alike. Part (b) is the benign extreme: when the features are independent given the label, whatever their association with the label and whatever the accuracy of the rule, a record contributes to the estimate only through which coordinates it reveals, and inclusion probabilities are all that a design controls. Any pair of inclusion vectors that a tree realizes is realized by a label-dependent mixture of forms, so adaptivity within the record has no first-order value once the label may be used. Without the label, a tree can use early answers as a proxy for the label and realize inclusion vectors that differ between the classes, which no label-blind mixture of forms does; the factor two in (ii) bounds this advantage, and it is nearly attained: if one feature is a 1\%-noisy copy of the label and the rule averages one of two groups of fair coins according to that feature, label-blind forms need $1.31$, $1.47$ and $1.57$ times the records of the tree that asks the copy first, at $b=3,4,5$, about $2(b-1)/b$ (\texttt{condind\_referee/test\_gap\_example.py}). The model is the one behind naive-Bayes classifiers and latent-class analyses with the outcome as the class; it is a submodel of $\mathcal M_2$, and (iii) shows that at balanced product laws the feature-by-feature interactions of $\mathcal M_2$ change nothing. Away from those laws they do: for random depth-$b$ rules at product laws with an unbalanced label, and at laws of $\mathcal M_{\rm ci}$ with label-dependent features, the pairwise constants computed by the solver differ by up to 15\% between forms and branching (\texttt{explore\_condind\_boundary.py}). Parts (c) and (d) say that, at laws without dependence,
the bound is far from attained when the model is parsimonious. Under the
order-$k$ model at uniform laws the price of pre-commitment is at most
$(b+1)^k$ for every contrast, and for the address contrast, which attains
$2^{b-1}$ without a model, it is exactly one under the pairwise model, a case of (b)(iii), and
$2^{b-1}$ again at $k=b+1$, when the model contains the interaction of all
$b$ queried items with the label. In between it is polynomial. Under the
order-$k$ model forms need at least $L_k(b)\ge\binom{b-1}{k-2}/b$ times the
records of the rule's own tree, a polynomial of degree $k-3$. Under the
three-way model this factor stays below four, but the rule's own tree is no
longer the best branching design once $b\ge6$: an interview that reads some
address bits and then several data bits consistent with them needs only
$U_3(b)=O(1/b)$, so the gap grows at least linearly in the cap ($U_3=1$ for
$b\le5$; the bound $L_3/U_3$ is $3.09$ at $b=6$, $4.10$ at $b=8$ and $9.28$ at
$b=20$). At product laws the exponential separation is thus a statement about
interactions of order growing with the cap. Part (e) shows that this reading
needs weak dependence and not only a parsimonious model. In the law of part
(e) the address bits, their one-hot copy and the product indicators are tied
by strong pairwise couplings, as dummy-coded categorical items and derived
indicators are in a real questionnaire. On the set that carries almost all of
the mass the target is a polynomial of degree two, so it is nearly in the
tangent space of the pairwise model, the model is of no help, and forms need $2^{b-1}/b$ times the records
of the rule's tree, within a factor $b$ of the bound of part (a). The order
of a log-linear model therefore measures parsimony only relative to the
coding of the features. The values of the constants along a path from the
uniform law to this law, and at pairwise fits of noisy address populations,
are in Table~\ref{tab:dependence}. The exact worst-case exponent at product laws,
between $\max\{k-3,1\}$ and $k$ for $k\ge3$, is open, and so is a bound of
type (c) at laws with dependence among the features, where conditional
expectations of low-order monomials are no longer low-order; Appendix~\ref{app:modelexp}
computes the constants there.

\begin{proof}
\emph{(a)} Fix $S\in\mathcal S_\tau$ and let $G_S=\{x:S_\tau(x)=S\}$. The
event $G_S$ is determined by $x_S$: starting at the root, follow $\tau$ on
the coordinates in $S$; the path of $x$ stays inside $S$ and is padded to $S$
exactly when $S_\tau(x)=S$. On $G_S$ the transcript $\ell_\tau(x)$ is a
function of $x_S$. Hence, by Jensen's inequality for the conditional
expectation given the coarser information,
\[
 J_S(h)\ \ge\ \E\big[\E[h\mid X_S,Y]^2\mathbf 1_{G_S}\big]\ \ge\ \E\big[\E[h\mid\ell_\tau(X),Y]^2\mathbf 1_{G_S}\big],
\]
and summing over $S\in\mathcal S_\tau$, whose events $G_S$ partition $\X$,
gives $\sum_SJ_S(h)\ge J_\tau(h)$, which is \eqref{eq:infodom}. A tree that
is not padded, or repeats a question, is refined by a padded tree, whose
information $J$ is at least as large by the same Jensen step, so an optimal
branching design may be taken to consist of padded trees. If $q$ is a
design on trees that attains $W^{\mathcal T}_{\rm branch}(P)$, the fixed
design $\sum_\tau q_\tau\,\mathrm{Unif}(\mathcal S_\tau)$ has
$J(h)\ge J_q(h)/\max_\tau|\mathcal S_\tau|$ for every $h$, so the supremum in
\eqref{eq:WT} grows by at most $\max_\tau|\mathcal S_\tau|\le K_{m,b}$
(Appendix~\ref{app:universal}). For $r$ rounds, replace the forms by the
$r$-round interviews of the proof of Theorem~\ref{thm:roundsfull}(a), which
follow $\tau$ for $r-1$ rounds and ask one of at most
$\min\{2^{b-r},\binom{m-r}{b-r}\}$ guessed completions in the last batch. For every profile exactly one
completion is correct, the event that a given completion is correct is
determined by that interview's transcript, on it the transcript determines
$\ell_\tau$, and the completions are drawn uniformly, so the two inequalities
above apply with these events in place of the $G_S$.

\emph{(b)} The scores of $\mathcal M_{\rm ci}$ at $P$ are the functions
$h(x,y)=c(y)+\sum_jg_j(x_j,y)$ with $\E c(Y)=0$ and $\E[g_j(X_j,Y)\mid Y]=0$,
that is, $g_j(x_j,y)=\beta_{y,j}\phi^y_j(x_j)$; these three kinds of terms are
orthogonal in $L^2(P)$. Let $\iota$ be a padded interview of any class, chosen
possibly as a function of $y$, and let $S(x)$ be the set of coordinates it
queries on the profile $x$. A transcript is a cylinder event $\{X_S=u\}$, and
given $Y=y$ the coordinates outside $S$ keep their conditional laws, so
\[
 \E[h\mid\ell_\iota(X),Y]=c(Y)+\sum_{j\in S(X)}g_j(X_j,Y).
\]
The event that $\iota$ queries $j$ is determined by the answers to the
coordinates queried before $j$ (in earlier batches, for an $r$-round
interview), which are independent of $X_j$ given $Y$. Hence
$\E[g_j\mathbf 1\{j\in S(X)\}\mid Y]=0$ and
$\E[g_j^2\mathbf 1\{j\in S(X)\}\mid Y=y]=\pi^\iota_j(y)\beta_{y,j}^2$ with
$\pi^\iota_j(y)=P(j\in S(X)\mid Y=y)$. For $i\ne j$, split the event
$\{i,j\in S(X)\}$ according to which of the two is queried later (in an arbitrary order if they are in the same batch), say $j$ at
a node $v$ of the interview: the event of reaching $v$ and the value $g_i$
depend on coordinates other than $j$, so
$\E[g_ig_j\mathbf 1\{\text{$j$ is queried at $v$ after $i$}\}\mid Y]=0$. Therefore
\[
 J_\iota(h)=\E[c(Y)^2]+\sum_yP(y)\sum_j\pi^\iota_j(y)\,\beta_{y,j}^2 ,
\]
and for a design the same holds with $\pi_j(y)=\sum_\iota q_\iota(y)\pi^\iota_j(y)$.
The projection of $t-\Delta$ on $\mathcal T_{\rm ci}$ has the components
$c^\ast(y)=\E[t\mid Y=y]-\Delta$ and
$\beta^\ast_{y,j}=\E[t\phi^y_j\mid Y=y]=\pm\hat d^{\,y}_j$, so by the
Cauchy--Schwarz inequality in these orthogonal coordinates the supremum in
\eqref{eq:WT} equals
\[
 \Var_P(\E[t\mid Y])+\sum_yP(y)\sum_j\frac{(\hat d^{\,y}_j)^2}{\pi_j(y)} .
\]
A design enters only through the inclusion vectors $\pi(y)$. Every padded
interview asks exactly $b$ coordinates of every record, so each $\pi(y)$ lies
in the hypersimplex $H_b=\{\pi\in[0,1]^m:\sum_j\pi_j=b\}$, whose extreme points
are the indicator vectors of the forms; the objective is nonincreasing in $\pi$, so the constraint $\sum_j\pi_j\le b$ of $M_b$ may be taken with equality. Label-dependent mixtures of forms
therefore realize every pair $(\pi(-1),\pi(1))\in H_b^2$, which is all that
any design of any class can realize; minimizing over each $\pi(y)$ separately
gives (i). Without the label, mixtures of forms realize exactly the pairs
with $\pi(-1)=\pi(1)$, which gives the formula for $W^{\rm ci}_{\rm fixed}$,
and the other classes lie between it and (i) because the menus are nested.
If $\pi^\ast(y)$ are optimal in (i), the common vector
$\pi=(\pi^\ast(-1)+\pi^\ast(1))/2\ge\pi^\ast(y)/2$ at most doubles every
term, which gives the last inequality of (ii). The form of the minimizer of
$M_b$ follows from the Lagrange conditions of the convex program. For (iii),
at a product law with a balanced label the standardized monomials
$\phi_A$, $A\subseteq[m+1]$, are orthonormal and $\langle t,\phi_A\rangle=0$
unless $A$ contains the label, because $\E Y=0$; the gradient in
$\mathcal T_2$ is supported on $\{m+1\}$ and the pairs $\{j,m+1\}$, where it
equals $\E d$ and $\hat d_j$. For an interview chosen without the label, the
information matrix in the basis $\{\phi_A\}$ has no entries between sets with
and without the label, since $\E\phi_{m+1}=0$ and $\phi_{m+1}^2$ has mean one,
and on the sets with the label it is the diagonal matrix with entries one and
$\pi^\iota_j$ by the computation above. The value $v^{\mathsf T}I^\dagger v$
is then $(\E d)^2+\sum_j\hat d_j^{\,2}/\pi_j$ for every design, and the
inclusion vectors of all classes fill $H_b$.

\emph{(c)} Let $\phi_j$ be the standardized $j$th variable and
$\phi_A=\prod_{j\in A}\phi_j$; at a product law the $\phi_A$,
$A\subseteq[m+1]$, are orthonormal, and $\{\phi_A:1\le|A|\le k\}$ is a basis of
$\mathcal T_k$. Write $h=\sum_Aa_A\phi_A$. A transcript of a padded interview
is a cylinder event $\{X_S=u\}$, so the coordinates outside $S$ keep their
marginal laws given the transcript and
$\E[\phi_A\mid\ell,Y]=\phi_A\mathbf 1\{A\subseteq S\cup\{m+1\}\}$. For a form
this gives $J_S(h)=\sum_{A\subseteq S\cup\{m+1\}}a_A^2$. For the tree,
\[
 J_\tau(h)=\E\Big[\Big(\sum_{A\in\mathcal A(X)}a_A\phi_A\Big)^2\Big]
   \le\E\Big[\Big(\sum_{A\in\mathcal A(X)}\phi_A^2\Big)\Big(\sum_{A\in\mathcal A(X)}a_A^2\Big)\Big]
   \le\nu_k\sum_Aa_A^2\,P\big(A\setminus\{m+1\}\subseteq S_\tau(X)\big),
\]
where $\mathcal A(x)$ is the collection of sets $A\subseteq S_\tau(x)\cup\{m+1\}$ with $1\le|A|\le k$,
because $\phi_j^2\le(1-\lambda)/\lambda$ and at most $\binom{b+1}{j}$ sets of
size $j$ lie in $S_\tau(X)\cup\{m+1\}$. The last sum is
$\sum_Sq^\tau_SJ_S(h)$. Mixing over the trees of an optimal branching design
gives the bound with $\nu_k$, and (a) gives the bound with $K_{m,b}$.

\emph{(d)} Here $\Delta=0$, $t=Yd$ and $\phi_A=\chi_A$. The gradient is
$v_A=\E[t\chi_A]$, which vanishes unless $A=B\cup\{j\}\cup\{m+1\}$ with $B$ a
set of address bits and $j$ a data bit, and then $|v_A|=2^{-\kappa}$ since
$\hat d(B\cup\{j\})=\E[\mathbf 1\{\text{address}=j\}\chi_B]=\pm2^{-\kappa}$. In
$\mathcal T_k$ there are $D\,N(\kappa,k-2)$ such sets. By the computation in
(c), the information matrix of a fixed design $q$ in the basis $\chi_A$ is
diagonal with entries $c_A(q)=P_q(A\setminus\{m+1\}\subseteq S)$, so
\[
 W^{(k)}_{\rm fixed}(P)=\min_q\sum_A\frac{v_A^2}{c_A(q)}\ \ge\ \min_q\frac{(\sum_A|v_A|)^2}{\sum_Ac_A(q)}
   \ \ge\ \frac{N(\kappa,k-2)^2}{\max_{s}(b-s)N(s,k-2)},
\]
by Cauchy--Schwarz, since a form with $s$ address bits and $b-s$ data bits
contains $(b-s)N(s,k-2)$ of the sets $B\cup\{j\}$. The maximum is at most
$bN(\kappa,k-2)$, which gives $L_k(b)\ge N(\kappa,k-2)/b\ge\binom{b-1}{k-2}/b$.
For $k=3$, $N(s,1)=1+s$ and $\max_s(b-s)(1+s)=\lfloor(b+1)^2/4\rfloor$; for
$k=b+1$, $N(s,\kappa)=2^s$ and the maximum is $2^\kappa$ at $s=\kappa$. The
tree that reads the address and the selected bit observes $t$, so
$W^{(k)}_{\rm branch}\le W_{\rm branch}=\sigma_P^2=1$. For the upper bound at
$k=3$, draw $s=\lfloor b/2\rfloor$ address bits and $b-s$ data bits uniformly.
Then $c_A=(b-s)/D$ for the $D$ sets with $B=\emptyset$ and
$c_A=s(b-s)/(\kappa D)$ for the $\kappa D$ sets with $|B|=1$, so
\[
 W^{(3)}_{\rm fixed}(P)\le\frac1{b-s}+\frac{\kappa^2}{s(b-s)}
   \le\frac2b+\frac{4(b-1)^2}{b^2-1}=\frac2b+4-\frac8{b+1}<4,
\]
using $s(b-s)\ge(b^2-1)/4$ and $b-s\ge b/2$. At $b=3$ and $b=4$ the first
bound equals $2.50$ and $2.75$, the values of Table~\ref{tab:addressmodel}.

For the bound on branching at $k=3$ we use weak duality in \eqref{eq:WT}: if
a design $q$ has scores $H_\iota(\ell_\iota,Y)$ with
$\sum_\iota q_\iota\E[H_\iota h]=\langle t-\Delta,h\rangle_P$ for all
$h\in\mathcal T$, then
$\langle t-\Delta,h\rangle_P^2=(\sum_\iota q_\iota\E[H_\iota\E[h\mid\ell_\iota,Y]])^2\le\sum_\iota q_\iota\E[H_\iota^2]\,J_q(h)$
by Cauchy--Schwarz, so $W^{\mathcal T}_\A(P)\le\sum_\iota q_\iota\E[H_\iota^2]$.
Fix $s,r\ge1$ with $s+r=b$ and $r\le2^{\kappa-s}$. With probability
$1-\epsilon$, draw a set $S_a$ of $s$ address bits uniformly, read them, and
then read a set $R$ of $r$ data bits drawn uniformly among the $2^{\kappa-s}$
data bits whose index agrees with the answers on $S_a$; the score is
$c\,Y\sum_{j\in R}Z_j/(1-\epsilon)$ with $c=\kappa/(rs)$. With probability
$\epsilon$, read a uniform set $R'$ of $b$ data bits; the score is
$-\delta\,Y\sum_{j\in R'}Z_j/\epsilon$ with $\delta=(\kappa-s)/(sb)$. This is a
mixture of two-round interviews. Given the profile, data bit $j$ belongs to
$R$ with probability $r2^{s-\kappa}\mathbf 1\{j\text{ agrees with the address on }S_a\}$,
and averaging the indicator over $S_a$ gives
$2^{-s}\sum_BP(B\subseteq S_a)\chi_B(\text{address})\chi_B(j)$ over sets $B$ of
address bits, with $P(B\subseteq S_a)=1$ for $B=\emptyset$ and $s/\kappa$ for
$|B|=1$. The conditional mean of the score given $(X,Y)$ is therefore
$Y\sum_jZ_j2^{-\kappa}\sum_Bw_B\chi_B(\text{address})\chi_B(j)$ with
$w_\emptyset=cr-\delta b=1$ and $w_B=crs/\kappa=1$ for $|B|=1$, whereas
$t=Y\sum_jZ_j2^{-\kappa}\sum_B\chi_B(\text{address})\chi_B(j)$. The two
functions have the same coefficients on all sets of at most three
coordinates, hence the same inner product with every $h\in\mathcal T_3$, and
\[
 W^{(3)}_{\rm branch}(P)\le\frac{c^2r}{1-\epsilon}+\frac{\delta^2b}{\epsilon},
 \qquad\min_\epsilon=\big(c\sqrt r+\delta\sqrt b\big)^2
 =\Big(\frac{\kappa}{s\sqrt r}+\frac{\kappa-s}{s\sqrt b}\Big)^2 ,
\]
since the data bits are independent of the address and of each other. The
choice $s=\kappa$, $r=1$ is the rule's own tree and gives one. With
$s=\lceil3b/4\rceil$, which is admissible for large $b$, the two terms are
$(8/3+o(1))/\sqrt b$ and $(1/3+o(1))/\sqrt b$, so $U_3(b)=(9+o(1))/b$, and
$L_3(b)\to4$.

For $k=2$ the gradient is supported on the $D$ pairs $\{j,m+1\}$ with $j$ a
data bit. Forms of $b$ data bits drawn uniformly have $c_{\{j,m+1\}}=b/D$ and
value $D\,2^{-2\kappa}D/b=1/b$, and the display gives
$W^{(2)}_{\rm fixed}\ge D^2\,2^{-2\kappa}/b=1/b$. For a branching design, let
$I$ be its information matrix and $I_{vv}$ the block of the $D$ pairs. Then
$v^{\mathsf T}I^\dagger v=\sup_u(2v^{\mathsf T}u-u^{\mathsf T}Iu)\ge\sup_{u\text{ supported on the block}}(2v^{\mathsf T}u-u^{\mathsf T}I_{vv}u)=v^{\mathsf T}I_{vv}^{\dagger}v$,
with the value $+\infty$ if some $c_j$ below vanishes.
For data bits $j\ne l$ the entry of a tree is
$\E[X_jX_l\mathbf 1\{j,l\in S_\tau(X)\}]$, which is zero: on the event that $l$
is read after $j$, the indicator is determined by the answers read before
$l$, and $X_l$ is independent of them with mean zero; the other order is
symmetric. So $I_{vv}$ is diagonal with entries $c_j=P(j\text{ is read})$,
$\sum_jc_j\le b$, and $v^{\mathsf T}I_{vv}^{\dagger}v=\sum_j2^{-2\kappa}/c_j\ge1/b$.
Mixtures of trees have the same diagonal block structure by linearity.

\emph{(c), stability.} Let $k\le b$, let $h_a=\sum_Aa_A(\chi_A-\E_P\chi_A)$
for a coefficient vector $a$ of Euclidean norm one over the sets with
$1\le|A|\le k$, and write $J^P$ for information at the law $P$. The menus of
forms and of padded depth-$b$ trees are finite, and for full-support laws
$(P,a)\mapsto J^P_\iota(h_a)$ is continuous, so
$\omega(\delta)=\sup\{|J^P_\iota(h_a)-J^{P_0}_\iota(h_a)|:\|P-P_0\|\le\delta,\ \|a\|=1,\ \iota\}$
tends to zero with $\delta$. Every set $A$ with $|A|\le k\le b$ lies in
$S\cup\{m+1\}$ for some form $S$, so the uniform design $u$ on all forms has
$c_0=\min_{\|a\|=1}J^{P_0}_u(h_a)>0$. For a tree $\tau$ and $\eta\in(0,1)$ put
$q^{\tau,\eta}=(1-\eta)q^\tau+\eta u$ with $q^\tau$ from (c) at $P_0$. Then
$J^{P_0}_{q^{\tau,\eta}}(h_a)\ge\max\{(1-\eta)J^{P_0}_\tau(h_a)/\nu_k,\ \eta c_0\}$,
and for $\|P-P_0\|\le\delta$ with $\omega(\delta)\le\eta c_0/2$,
\[
 J^P_\tau(h_a)\le J^{P_0}_\tau(h_a)+\omega\le\frac{\nu_k}{1-\eta}\big(J^P_{q^{\tau,\eta}}(h_a)+\omega\big)+\omega
 \le\frac{\nu_k}{1-\eta}\Big(1+\frac{4\omega}{\eta c_0}\Big)J^P_{q^{\tau,\eta}}(h_a),
\]
using $J^P_{q^{\tau,\eta}}(h_a)\ge\eta c_0/2$ and $\nu_k\ge1$. The tangent
space $\mathcal T_k$ at $P$ consists of the multiples of the $h_a$, so the
fixed design $\sum_\tau q_\tau q^{\tau,\eta}$ built from an optimal branching
design $q$ at $P$ has information at least
$(1-\eta)(1+4\omega/(\eta c_0))^{-1}/\nu_k$ times that of $q$ in every
direction of $\mathcal T_k$. Choosing $\eta$ and then $\delta$ small gives
the factor $(1+\epsilon)\nu_k$. The radius $\delta$ depends on $m$, $b$, $k$,
$\lambda$ and $\epsilon$ and is not explicit.

\emph{(e)} Write the address as $a\in\{-1,1\}^\kappa$, identify a data index
$j$ with its sign vector, and let $\langle a,j\rangle\in\{-\kappa,\ldots,\kappa\}$
be their inner product, which equals $\kappa$ exactly when $j$ is the index
selected by $a$. With $z_j=(1+Z_j)/2$ define the quadratic function
\[
 \mathcal Q(a,s,z,p)=\sum_js_j\langle a,j\rangle-(\kappa+1)\Big(\sum_js_j-1\Big)^2
   -\sum_j\big(3p_j+s_jz_j-2p_js_j-2p_jz_j\big),
\]
and let $P_\beta$ be the law with $P_\beta(x)\propto e^{\beta\mathcal Q(x)}$ and an
independent uniform label. $\mathcal Q$ is a polynomial of degree two in the binary
features, so $P_\beta\in\mathcal M_2$ for every $\beta\ge0$, with full
support. The last sum is the standard penalty for a product of two binary
variables \citep{boros2002}: each term is nonnegative and vanishes exactly
when $p_j=s_jz_j$. If exactly one $s_j$ equals one, the first two terms equal
$\langle a,j\rangle\le\kappa$, with equality exactly when $j$ is the selected
index; if none does they equal $-(\kappa+1)$; and if $n\ge2$ do, they are at
most $n\kappa-(\kappa+1)(n-1)^2<\kappa$. So $\mathcal Q\le\kappa$ with equality
exactly on the set where $s$ is the one-hot code of the address and
$p_j=s_jz_j$ for all $j$, and as $\beta\to\infty$ the law $P_\beta$ converges
to the law $P_\infty$ under which the address, the data bits and the label
are independent and uniform and $s$ and $p$ are these functions of them.

On the support of $P_\infty$,
$d=Z_{a}=\sum_js_j(2z_j-1)=2\sum_jp_j-1$, so $t=h_0$ almost surely for the
polynomial $h_0=Y(2\sum_jp_j-1)$ of degree two in $(p,Y)$, whose centered
version lies in $\mathcal T_2\subseteq\mathcal T_k$ at every $P_\beta$. By
\eqref{eq:WT},
\[
 W^{(k)}_{\rm fixed}(P_\beta)\ \ge\ \frac{\mathrm{Cov}_{P_\beta}(t,h_0)^2}{\max_SJ^{P_\beta}_S(h_0-\E_{P_\beta}h_0)} .
\]
For a bounded function $h$ and an interview $\iota$,
$J^P_\iota(h)=\sum_c\E_P[h\mathbf 1_c]^2/P(c)$ over the transcript cells, and
a cell contributes at most $\|h\|_\infty^2P(c)$, so $J^{P_\beta}_S\to
J^{P_\infty}_S$ and the right-hand side converges to
$1/\max_SJ^{P_\infty}_S(t)$, because $\mathrm{Var}_{P_\infty}(t)=1$. Fix a
form $S$ of $b$ of the $m$ features, with data bits $R_z$ and product bits
$R_p$. Under $P_\infty$ and given that the selected index is $j$, the
features in $S$ other than $Z_j$ and $p_j$ are functions of the address and
of data bits independent of $Z_j$, while $p_j=z_j$; hence the conditional
mean of $Z_j$ given $X_S$ and the selected index $j$ is $Z_j$ or $2p_j-1$ if
$j\in R_z\cup R_p$ and zero otherwise. Therefore
$|\E[d\mid X_S]|\le P(\text{selected index}\in R_z\cup R_p\mid X_S)$ and
\[
 J^{P_\infty}_S(t)=\E\big[\E[d\mid X_S]^2\big]\le P(\text{selected index}\in R_z\cup R_p)\le\frac bD .
\]
So $\liminf_\beta W^{(k)}_{\rm fixed}(P_\beta)\ge D/b$. The tree that reads the
address and the selected data bit computes $d$ at every profile, so
$W^{(k)}_{\rm branch}(P_\beta)\le W_{\rm branch}(P_\beta)=\mathrm{Var}_{P_\beta}(t)=1$,
because $t=\pm1$ and the independent uniform label gives $\Delta(P_\beta)=0$.
Taking $\beta$ large gives the claim, and $P_\beta$ puts mass tending to one
on the support of $P_\infty$. For $b=2$ the bound $D/b$ equals one.
\end{proof}

\paragraph{Computed constants for the address contrast.}
Table~\ref{tab:addressmodel} lists the constants of
Theorem~\ref{thm:structuredfull}(d) for $b=3$ and $b=4$, computed by the
model version of the solver (Appendix~\ref{app:modelexp}). Under the
three-way model the gap is $2.50$ and $2.75$ against the lower bounds
$L_3(3)=2.25$ and $L_3(4)=2.67$, and at these caps the optimal branching
design is the rule's own tree ($U_3=1$ for $b\le5$); under the pairwise model
both classes need $1/b$ of the records that the design-based branching
optimum needs.
\begin{table}[h]
\centering
\caption{The address contrast at the uniform law within the order-$k$ log-linear model: constants relative to $\sigma_P^2=1$, computed by the solver (certified gaps below $10^{-5}$). Complete: bound for complete records. $L_k(b)$: the lower bound of Theorem~\ref{thm:structuredfull}(d) on the fixed constant. Without a model the ratio is $2^{b-1}$.}
\label{tab:addressmodel}
\small
\begin{tabular}{rrrrrrrr}
\toprule
$b$ & $m$ & $k$ & Complete & Fixed & Branch & Ratio & $L_k(b)$ \\
\midrule
3 & 6 & 2 & 0.250 & 0.333 & 0.333 & 1.00 & 0.333 \\
3 & 6 & 3 & 0.750 & 2.500 & 1.000 & 2.50 & 2.250 \\
3 & 6 & 4 & 1.000 & 4.000 & 1.000 & 4.00 & 4.000 \\
4 & 11 & 2 & 0.125 & 0.250 & 0.250 & 1.00 & 0.250 \\
4 & 11 & 3 & 0.500 & 2.750 & 1.000 & 2.75 & 2.667 \\
\bottomrule
\end{tabular}
\end{table}

